\documentclass[11pt]{article}

\usepackage[T1]{fontenc}
\usepackage[utf8]{inputenc}
\usepackage{lmodern}
\usepackage{amsmath,amssymb,amsthm,mathtools}
\usepackage{booktabs}
\usepackage{enumitem}
\usepackage[margin=1in]{geometry}
\usepackage[colorlinks=true,linkcolor=blue,citecolor=blue,urlcolor=blue]{hyperref}
\usepackage{microtype}

\newtheorem{theorem}{Theorem}[section]
\newtheorem{lemma}[theorem]{Lemma}
\newtheorem{proposition}[theorem]{Proposition}
\newtheorem{corollary}[theorem]{Corollary}
\theoremstyle{remark}
\newtheorem{remark}[theorem]{Remark}

\newcommand{\RR}{\mathbb{R}}
\newcommand{\CC}{\mathbb{C}}
\newcommand{\HH}{\mathbb{H}}
\newcommand{\fg}{\mathfrak{g}}
\newcommand{\fh}{\mathfrak{h}}
\newcommand{\fs}{\mathfrak{s}}
\newcommand{\fz}{\mathfrak{z}}
\newcommand{\so}{\mathfrak{so}}
\newcommand{\su}{\mathfrak{su}}
\newcommand{\spalg}{\mathfrak{sp}}
\newcommand{\End}{\operatorname{End}}
\newcommand{\rank}{\operatorname{rank}}
\newcommand{\Lie}{\operatorname{Lie}}
\newcommand{\Iso}{\operatorname{Iso}}
\newcommand{\Spin}{\operatorname{Spin}}
\newcommand{\SU}{\operatorname{SU}}
\newcommand{\SO}{\operatorname{SO}}
\newcommand{\U}{\operatorname{U}}
\newcommand{\Sp}{\operatorname{Sp}}
\newcommand{\Gtwo}{G_2}
\newcommand{\Ttri}[1]{\binom{#1}{2}}

\title{A Small-Module Lemma for the Skew-Hermitian Dimension Problem}
\author{Draft appendix}
\date{June 2026}

\begin{document}
\maketitle

\begin{abstract}
This note proves the representation-theoretic small-module lemma used in the triangular-deficit approach to the finite-dimensional skew-Hermitian Lie algebra problem.  The conclusion is that an irreducible compact connected orthogonal representation whose Lie-algebra dimension has triangular deficit must be one of the standard sphere-transitive atoms, namely the standard real representations of \(\SU(k)\) or \(\U(k)\) on \(\CC^k\), or the standard \(\Gtwo\)-module on \(\RR^7\).  Full orthogonal blocks occur only on the equality boundary and not in the strict deficit case.
\end{abstract}

\section{Triangular support}

For \(r\geq 2\), put
\[
  T_r=\binom r2.
\]
Define the triangular support function
\[
  w(N)=\min\left\{\sum_i r_i:
       N=\sum_i \binom{r_i}{2},\ r_i\geq 2\right\},
       \qquad w(0)=0.
\]
Since \(T_2=1\), the function is defined for every non-negative integer \(N\).
It is useful to keep in mind that
\[
  N\in P(n):=\left\{\sum_i \binom{r_i}{2}:\sum_i r_i\leq n,\ r_i\geq 2\right\}
  \quad\Longleftrightarrow\quad
  w(N)\leq n.
\]
Thus the obstruction wanted in the Banach-space problem is exactly the inequality
\(w(\dim Z(X))\leq \dim X\).

\begin{lemma}[Elementary estimates for \(w\)]\label{lem:w-estimates}
The following estimates hold.
\begin{enumerate}[label=\textup{(\roman*)}]
\item \(w(A+B)\leq w(A)+w(B)\).
\item \(w(T_r)=r\).
\item For \(k\geq 3\),
\[
   w(k^2-1)\leq 4k-4.
\]
\item For \(k\geq 2\),
\[
   w(k^2)\leq 2k+1.
\]
\item For \(k\geq 1\),
\[
   w(k(2k+1))\leq 2k+1.
\]
\end{enumerate}
\end{lemma}

\begin{proof}
Subadditivity follows by concatenating triangular decompositions.  The identity \(w(T_r)=r\) is immediate from the one-block decomposition, and no smaller support can contain \(T_r\) dimensions of \(\so(r)\).
For (iii), use
\[
  k^2-1=\binom{k+1}{2}+\binom{k-1}{2}+(k-2)\binom22,
\]
which has support \((k+1)+(k-1)+2(k-2)=4k-4\).  For (iv), use
\[
  k^2=\binom{k+1}{2}+\binom k2,
\]
which has support \(2k+1\).  Finally,
\[
  k(2k+1)=\binom{2k+1}{2},
\]
so \(w(k(2k+1))\leq 2k+1\).
\end{proof}

We shall also use the small explicit values
\[
\begin{gathered}
  w(1)=2,
  \quad w(3)=3,
  \quad w(4)=5,
  \quad w(6)=4,
  \quad w(8)=8,
  \quad w(9)=7,\\
  w(14)=10,
  \quad w(15)=6,
  \quad w(16)=8.
\end{gathered}
\]

\section{Statement of the lemma}

\begin{theorem}[Small-Module Lemma, strict form]\label{thm:small-module}
Let \(H\leq \SO(E)\) be compact, connected, effective, and irreducible over \(\RR\).  Put
\[
  d=\dim_{\RR}E,
  \qquad
  h=\dim H.
\]
If
\[
  w(h)>d,
\]
then, up to finite covers, dualizing complex modules, and the standard low-rank isomorphisms, the representation is one of the following:
\[
  \SU(k)\curvearrowright \CC^k_{\RR},\qquad k\geq 3,
\]
\[
  \U(k)\curvearrowright \CC^k_{\RR},\qquad k\geq 2,
\]
or
\[
  \Gtwo\curvearrowright \RR^7.
\]
Every representation in this list is transitive on the Euclidean unit sphere of \(E\).  Conversely, not every value of \(k\) in the displayed \(\SU(k)\)- and \(\U(k)\)-families satisfies the strict inequality; the theorem says only that these are the only possible strict-deficit shapes.
\end{theorem}

\begin{remark}[Equality blocks]
The full orthogonal block
\[
  \SO(d)\curvearrowright \RR^d
\]
has
\[
  \dim \SO(d)=\binom d2,
  \qquad
  w\!\left(\binom d2\right)=d.
\]
Thus a full orthogonal block is an equality case, not a strict-deficit case.  This is why the strict lemma above does not list \(\SO(d)\curvearrowright\RR^d\), although such blocks are harmless in the non-strict bookkeeping of the general proof.
\end{remark}

\section{Reduction to compact reductive representation theory}

Let \(\fh=\Lie(H)\).  Since \(H\) is compact,
\[
  \fh=\fz\oplus \fs,
  \qquad
  \fs=[\fh,\fh]=\fs_1\oplus\cdots\oplus\fs_r,
\]
where \(\fz\) is abelian and each \(\fs_i\) is compact simple.

\begin{lemma}[The center has dimension at most one]\label{lem:center-one}
If \(E\) is irreducible over \(\RR\), then \(\dim \fz\leq 1\).
\end{lemma}

\begin{proof}
By Schur's lemma, the real commuting algebra
\[
  D:=\End_{\fh}(E)
\]
is one of \(\RR,\CC,\HH\).  The center \(\fz\) acts by skew-symmetric elements of \(D\).  In \(\RR\) there are no nonzero skew elements.  In \(\CC\), the skew part is \(\RR i\).  In \(\HH\), the skew part is \(\operatorname{Im}\HH\), whose abelian subalgebras are one-dimensional.  Since \(\fz\) is abelian and the action is effective on the Lie algebra, \(\dim\fz\leq 1\).
\end{proof}

The case \(\fs=0\) is harmless.  Indeed, a nontrivial irreducible real representation of a connected abelian compact group is the usual circle action on \(\RR^2\); here \(h=1\), \(d=2\), and \(w(1)=2=d\).  Thus a strict deficit requires \(\fs\neq 0\).

\section{One simple factor}

The next proposition is the standard small-highest-weight input.  It is stated in the exact form needed here.

\begin{proposition}[Small irreducible modules for simple compact algebras]\label{prop:simple-table}
Let \(\fs\) be compact simple, and let \(U\) be a faithful irreducible real \(\fs\)-module.  Put \(s=\dim\fs\) and \(u=\dim_{\RR}U\).  If \(w(s)>u\), then either
\[
  (\fs,U)=(\su(k),\CC^k_{\RR})\quad (k\geq 3),
\]
up to dualizing the standard complex module, or
\[
  (\fs,U)=(\mathfrak g_2,\RR^7).
\]
The standard \(\so(u)\)-module satisfies equality \(w(\dim\so(u))=u\), not strict inequality.
\end{proposition}

\begin{proof}
We run through the simple types.

For \(\fs=\su(k)\), \(k\geq 3\), Lemma~\ref{lem:w-estimates} gives
\[
  w(\dim\su(k))=w(k^2-1)\leq 4k-4.
\]
The only irreducible real \(\su(k)\)-modules of real dimension below \(4k-4\) are the realifications of the standard module \(\CC^k\) and its dual, except for the low-rank isomorphism
\[
  \su(4)\cong\so(6),
  \qquad
  \Lambda^2\CC^4\cong\RR^6.
\]
This is the usual first-row consequence of the Weyl dimension formula for type \(A_{k-1}\): after the fundamental weights \(\omega_1,\omega_{k-1}\), the next smallest complex modules are \(\Lambda^2\CC^k\), \(\Lambda^{k-2}\CC^k\), and, in the small ranks, the adjoint module.  Their real dimensions are at least \(4k-4\), with the single real-type exception \(\Lambda^2\CC^4\), whose real dimension is \(6\).  For that exception,
\[
  \dim\su(4)=15=\binom62,
  \qquad
  w(15)=6,
\]
so it is an equality block.  Hence only the standard \(\su(k)\)-module can occur under strict deficit.

For \(\fs=\su(2)\), the nontrivial irreducible real modules have dimensions \(3,4,5,\ldots\).  Since \(\dim\su(2)=3\) and \(w(3)=3\), there is no strict deficit for \(\su(2)\) alone.

For \(\fs=\so(m)\), \(m\geq 5\),
\[
  \dim\so(m)=\binom m2,
  \qquad
  w(\dim\so(m))=m.
\]
The vector representation has dimension \(m\), hence equality.  All non-vector irreducible real modules have real dimension at least \(m\); the familiar spin and half-spin low-rank cases \(\Spin(5)\cong\Sp(2)\), \(\Spin(6)\cong\SU(4)\), \(\Spin(7)\), and triality for \(\Spin(8)\) still have real dimension at least \(m\), with equality only for vector-type blocks.  Thus no strict deficit occurs for type \(B\) or \(D\) outside the type \(A\) isomorphism already handled.

For \(\fs=\spalg(k)\),
\[
  \dim\spalg(k)=k(2k+1)=\binom{2k+1}{2},
  \qquad
  w(\dim\spalg(k))
  \leq 2k+1.
\]
The standard quaternionic module has real dimension \(4k\), and all other irreducible real modules are no smaller.  Since \(2k+1\leq 4k\), no strict deficit occurs.

For the exceptional algebras, the relevant numbers are as follows.
\[
\begin{array}{c|c|c|c}
\fs & \dim\fs & \text{upper bound for } w(\dim\fs) & \text{smallest real faithful module}\\
\midrule
\mathfrak g_2 & 14  & 10 & 7\\
\mathfrak f_4 & 52  & 16 & 26\\
\mathfrak e_6 & 78  & 13 & 54\\
\mathfrak e_7 & 133 & 23 & 56\\
\mathfrak e_8 & 248 & 32 & 248
\end{array}
\]
The displayed bounds for \(w\) are obtained, respectively, from
\[
  14=\binom52+\binom32+\binom22,
\]
\[
  52=\binom{10}{2}+\binom42+\binom22,
  \qquad
  78=\binom{13}{2},
\]
\[
  133=\binom{15}{2}+\binom82,
  \qquad
  248=\binom{22}{2}+\binom62+2\binom22.
\]
Thus \(\mathfrak f_4,\mathfrak e_6,\mathfrak e_7,\mathfrak e_8\) are safe.  For \(\mathfrak g_2\), the standard seven-dimensional module has
\[
  w(14)=10>7,
\]
and the next irreducible real module is the adjoint module of dimension \(14\), which is safe.  This leaves exactly \(\mathfrak g_2\curvearrowright\RR^7\).
\end{proof}

\begin{proposition}[Simple algebra plus one central circle]\label{prop:centered-simple}
Let \(\fh=\fs\oplus\RR\), where \(\fs\) is compact simple, and let \(E\) be a faithful irreducible real \(\fh\)-module.  If \(w(\dim\fh)>\dim_{\RR}E\), then
\[
  \fh\cong \su(k)\oplus\RR
\]
and \(E\cong\CC^k_{\RR}\) is the standard \(\U(k)\)-module, with \(k\geq 2\).  Equivalently, the strict-deficit centered simple atoms are contained in
\[
  \U(k)\curvearrowright \CC^k_{\RR}.
\]
\end{proposition}

\begin{proof}
A nonzero central circle gives a commuting complex structure on \(E\).  Therefore the underlying simple module must be of complex or quaternionic type.

If \(\fs=\su(k)\), then
\[
  \dim\fh=k^2,
  \qquad
  w(k^2)\leq 2k+1
\]
by Lemma~\ref{lem:w-estimates}.  The standard module \(\CC^k_{\RR}\) has real dimension \(2k\), and is therefore the only possible strict-deficit case.  Every nonstandard complex-type irreducible \(\su(k)\)-module has real dimension at least \(2k+2\), and usually much larger; hence \(w(k^2)\leq 2k+1\leq \dim_{\RR}E\) for those modules.  For \(k=2\), this is exactly the familiar atom
\[
  \U(2)=\Sp(1)\U(1)\curvearrowright \HH\cong\CC^2.
\]

If \(\fs=\spalg(k)\), the standard quaternionic module admits a central circle.  Here
\[
  \dim(\spalg(k)\oplus\RR)=\binom{2k+1}{2}+1,
\]
so
\[
  w(\dim(\spalg(k)\oplus\RR))
  \leq (2k+1)+w(1)
  =2k+3.
\]
The standard module has real dimension \(4k\).  For \(k\geq 2\), \(2k+3\leq 4k\), and for \(k=1\) this is the already listed \(\U(2)\)-module.  Nonstandard modules are larger.

If \(\fs=\so(m)\), the vector representation is real type, so it does not admit an effective central circle.  For the spin and half-spin modules that do admit a complex or quaternionic structure, the real dimension is at least \(m+2\); meanwhile
\[
  w\!\left(\binom m2+1\right)\leq m+2.
\]
The low-rank case \(\Spin(6)\cong\SU(4)\) gives the standard \(\U(4)\)-module, where \(w(16)=8=\dim_{\RR}\CC^4\), again equality rather than strict deficit.  The exceptional algebras remain safe after adding one to their dimensions; the upper bounds in Proposition~\ref{prop:simple-table} increase by at most \(w(1)=2\), still far below the smallest real module dimensions, except for \(\mathfrak g_2\), whose standard module is real type and so admits no effective central circle.
\end{proof}

\section{Several simple factors}

\begin{proposition}[Products are safe]\label{prop:products-safe}
Let
\[
  \fh=\fz\oplus\fs_1\oplus\cdots\oplus\fs_r,
  \qquad r\geq 2,
\]
be compact reductive with \(\dim\fz\leq 1\), and let \(E\) be a faithful irreducible real \(\fh\)-module.  Then
\[
  w(\dim\fh)\leq \dim_{\RR}E.
\]
In particular, no strict-deficit module has two or more simple factors.
\end{proposition}

\begin{proof}
Complexifying, an irreducible module for the semisimple part is a tensor product
\[
  W=W_1\otimes\cdots\otimes W_r
\]
of irreducible nontrivial complex modules for the simple factors, up to the usual passage from a complex module to either a real form or a realification.  Tensoring makes the module dimension grow multiplicatively, whereas the Lie-algebra dimension and the triangular support grow additively.  We spell out the only tight cases.

First suppose one factor is a possible type \(A\) atom, \(\su(k)\) on \(\CC^k\), \(k\geq 3\).  The only tensor partner that is genuinely tight is the quaternionic standard \(\su(2)\)-module: then the tensor has real dimension \(4k\), while the Lie-algebra dimension is \(k^2+2\), or \(k^2+3\) if the unique central circle acts effectively.  The triangular decompositions
\[
  k^2+2=\binom{k+2}{2}+\binom{k-1}{2},
  \qquad
  k^2+3=\binom{k+2}{2}+\binom{k-1}{2}+\binom22
\]
give
\[
  w(k^2+2)\leq 2k+1,
  \qquad
  w(k^2+3)\leq 2k+3\leq 4k.
\]
Thus this tight case is safe.  If the partner is the three-dimensional real \(\su(2)\)-module, the real tensor dimension is \(6k\), and the cruder estimate
\[
  (4k-4)+3+w(\dim\fz)
  \leq 4k+1
  \leq 6k
\]
is already enough.  Larger partners only increase the tensor dimension faster than the support.  If two type \(A\) standard factors \(\su(k)\) and \(\su(\ell)\), \(k,\ell\geq 3\), occur, then the real tensor dimension is at least \(2k\ell\), and
\[
  w(k^2-1)+w(\ell^2-1)+w(\dim\fz)
  \leq (4k-4)+(4\ell-4)+2
  \leq 2k\ell,
\]
because \((k-2)(\ell-2)\geq 0\).  The centered estimate is even easier if one of the factors uses the global central circle as a \(\U(k)\)-standard action, since \(w(k^2)\leq 2k+1\).

For the \(\mathfrak g_2\)-standard factor, the smallest possible tensor partner has real dimension \(3\).  Hence the real tensor dimension is at least \(21\), whereas
\[
  w(14)+w(3)+w(\dim\fz)
  \leq 10+3+2=15.
\]
Thus a \(\mathfrak g_2\)-factor is safe once it is coupled to another simple factor.

It remains to check the quaternionic-quaternionic tensors, since these are the only products for which the real form of the tensor can have dimension half of the naive product of the underlying real dimensions.  The tight family is
\[
  \Sp(p)\times\Sp(q)
\]
acting on the real form of the tensor product of the two standard quaternionic modules.  Its real dimension is \(4pq\), and
\[
  \dim\spalg(p)+\dim\spalg(q)
  =\binom{2p+1}{2}+\binom{2q+1}{2}.
\]
Therefore
\[
  w(\dim\spalg(p)+\dim\spalg(q))
  \leq (2p+1)+(2q+1)
  \leq 4pq
\]
for all \(p,q\geq 1\), except that \((p,q)=(1,1)\) gives equality in the sharper form
\[
  \Sp(1)\times\Sp(1)\cong\Spin(4)\curvearrowright\RR^4,
  \qquad
  w(6)=4.
\]
If a central circle acts effectively on the tensor, then the tensor has a commuting complex structure and its real dimension doubles in the tight real-type cases; the possible extra support \(w(1)=2\) is then harmless.

All remaining product tensors have at least one factor whose triangular support is bounded by its real module dimension, and no exceptional halving of the tensor dimension occurs.  Subadditivity of \(w\), together with the estimates in Propositions~\ref{prop:simple-table} and~\ref{prop:centered-simple}, then gives \(w(\dim\fh)\leq\dim_{\RR}E\).  Hence products never give strict deficit.
\end{proof}

\section{Proof of the Small-Module Lemma}

\begin{proof}[Proof of Theorem~\ref{thm:small-module}]
Let \(\fh=\fz\oplus\fs\) be the compact reductive Lie algebra of \(H\).  Lemma~\ref{lem:center-one} gives \(\dim\fz\leq 1\).  If \(\fs=0\), the representation is either trivial or the usual circle representation on \(\RR^2\), and there is no strict deficit.

If \(\fs\) has at least two simple factors, Proposition~\ref{prop:products-safe} gives
\[
  w(\dim\fh)\leq \dim_{\RR}E,
\]
contrary to the hypothesis.  Thus \(\fs\) is simple.

If \(\fz=0\), Proposition~\ref{prop:simple-table} leaves only the standard \(\SU(k)\)-module on \(\CC^k_{\RR}\), \(k\geq 3\), and the standard \(\Gtwo\)-module on \(\RR^7\).

If \(\dim\fz=1\), Proposition~\ref{prop:centered-simple} leaves only the standard \(\U(k)\)-module on \(\CC^k_{\RR}\), \(k\geq 2\).  These are precisely the displayed possibilities.

Finally, \(\SU(k)\) and \(\U(k)\) are transitive on the unit sphere in \(\CC^k\), and \(\Gtwo\) is transitive on the unit sphere in its standard seven-dimensional real module.  This proves the lemma.
\end{proof}

\section{Insertion into the deficit-summand proof}

The form of Theorem~\ref{thm:small-module} is exactly what the deficit-summand argument needs.  If an irreducible image \(H_i\leq\SO(E_i)\) satisfies
\[
  w(\dim H_i)>\dim E_i,
\]
then \(E_i\) carries one of the proper sphere-transitive atoms above.  In the Banach-space application, invariance of the norm under such a proper sphere-transitive group on a summand forces invariance under the full orthogonal group of that summand; hence the connected isometry group would enlarge unless the atom is coupled to another summand.

One technical correction is worth recording.  In the centralizer-coupling step, if \(Q\) is the image of the group on \(E\oplus F\) and \(A\) is its image on \(E\), write
\[
  0\longrightarrow K\longrightarrow Q\longrightarrow A\longrightarrow 0,
  \qquad
  \ell=\dim K.
\]
Then
\[
  \dim Q=\dim A+\ell,
\]
not merely \(\dim Q\leq \dim A+\ell\).  This equality is useful because \(w\) is not monotone in general.

\begin{thebibliography}{9}

\bibitem{BourbakiLie}
N. Bourbaki,
\emph{Lie Groups and Lie Algebras, Chapters 4--6},
Springer, 2002.

\bibitem{BorelTransitiveSpheres}
A. Borel,
\emph{Le plan projectif des octaves et les sph\`eres comme espaces homog\`enes},
C. R. Acad. Sci. Paris \textbf{230} (1950), 1378--1380.

\bibitem{KadetsMartinPaya}
V. Kadets, M. Mart\'in, and R. Pay\'a,
\emph{Recent progress and open questions on the numerical index of Banach spaces},
arXiv:math/0605781.

\bibitem{RosenthalLie}
H. P. Rosenthal,
\emph{The Lie algebra of a Banach space},
in: Banach spaces (Columbia, Mo., 1984), Lecture Notes in Math. \textbf{1166}, Springer, 1985, 129--157.

\end{thebibliography}

\end{document}
